\section{Dominance Theorems for Hardness Placement}\label{sec:implications}

Regime for this section: the mechanized Boolean-coordinate model [S\_bool] plus the architecture cost model defined below. The centralization-vs-distribution framing follows established software-architecture criteria about modular decomposition and long-run maintenance tradeoffs \cite{parnas1972criteria, brooks1987no, bass2012software, shaw1996software}.
Although phrased in software terms, the formal objects are substrate-agnostic: the same hardness-placement and amortization equations apply to physical architectures (for example, centralized vs distributed sensing/actuation computation) once mapped into the same coordinate/cost model.
\claimstamp{P\ref{prop:physical-claim-transport};T\ref{thm:physical-bridge-bundle},\ref{thm:centralization-dominance},\ref{thm:tax-conservation},\ref{thm:amortization}}{AR,LC}
\leanmeta{\LH{CT6},\LH{CT4},\LH{HD1},\LH{HD5},\LH{HD4}.}

In electrical implementations, the same architecture equations are energy equations: each distributed update path induces switching events and Joule expenditure, while centralized one-time structure pays setup cost and reduces repeated per-site electrical work. Under the declared thermodynamic lift, this is tracked as a typed transfer from bit-operation load to physical energy load.
\claimstamp{P\ref{prop:thermo-lift};C\ref{cor:neukart-vinokur}}{AR,H=cost-growth}
\leanmeta{\LH{CC76},\LH{CC77}.}

Equivalently, repeated distributed cycles create a thermal duty: more update cycles imply more bit operations, which imply more unavoidable dissipated heat and therefore larger heat-transfer/removal requirements in the same substrate class.
\claimstamp{P\ref{prop:decision-unit-time},\ref{prop:run-time-accounting},\ref{prop:thermo-lift}}{AR,H=cost-growth}
\leanmeta{\LH{DT13},\LH{DT16},\LH{CV17},\LH{CV18},\LH{CC76}.}

\subsection{Over-Specification as Diagnostic Signal}

\begin{corollary}[Persistent Over-Specification]\label{cor:overmodel-diagnostic-implication}
\claimstamp{C\ref{cor:overmodel-diagnostic-implication}}{H=succinct-hard}
In the mechanized Boolean-coordinate model, if a coordinate is relevant and omitted from a candidate set $I$, then $I$ is not sufficient.
\leanmeta{\LH{S2P6}}
\end{corollary}

\begin{proof}
This is the contrapositive of \LH{S2P3}.
\end{proof}

\begin{corollary}[Exact Relevance Identifiability Criterion]\label{cor:exact-identifiability}
\claimstamp{C\ref{cor:exact-identifiability}}{H=succinct-hard}
In the mechanized Boolean-coordinate model, for any candidate set $I$:
\[
I \text{ is exactly relevance-identifying}
\iff
\bigl(I \text{ is sufficient and } I \subseteq R_{\mathrm{rel}}\bigr),
\]
where $R_{\mathrm{rel}}$ is the full relevant-coordinate set.
\leanmeta{\LH{S2P1}}
\end{corollary}

\begin{proof}
This is exactly
\LH{S2P1},
with $R_{\mathrm{rel}}=\texttt{relevantFinset}$.
\end{proof}

\begin{remark}[Over-Specification Is Regime-Conditional, Not Default]\label{rem:overmodel-conditional}
In this paper's formal typing, persistent over-specification is admissible only under attempted competence failure: after an attempted exact procedure, exact relevance competence remains unavailable in the active regime, so integrity retains uncertified coordinates (Section~\ref{sec:interpretive-foundations}; Section~\ref{sec:engineering-justification}).\claimstamp{P\ref{prop:attempted-competence-matrix}}{AR}
Once exact competence is available in the active regime (Corollaries~\ref{cor:practice-bounded}--\ref{cor:practice-symmetry} together with Corollary~\ref{cor:exact-identifiability}), certifiably irrelevant coordinates are removable; persistent over-specification is irrational for the same objective (verified total-cost minimization).\claimstamp{T\ref{thm:tractable};P\ref{prop:attempted-competence-matrix}}{TA}
\leanmeta{\LH{IC18},
\LH{IC25},
\LH{IC7},
\LH{CC81},
\LH{CC82},
\LH{CC84},
\LH{CC71},
\LH{CC73},
\LH{CC72},
\LH{S2P1}.}
\end{remark}

\subsection{Architectural Decision Quotient}

\begin{definition}[Architectural Decision Quotient]
For a software system with configuration space $S$ and behavior space $B$:
\[
\text{ADQ}(I) = \frac{|\{b \in B : b \text{ achievable with some } s \text{ where } s_I \text{ fixed}\}|}{|B|}
\]
\end{definition}

\begin{proposition}[ADQ Ordering]\label{prop:adq-ordering}
For coordinate sets $I,J$ in the same system, if $\mathrm{ADQ}(I) < \mathrm{ADQ}(J)$, then fixing $I$ leaves a strictly smaller achievable-behavior set than fixing $J$.
\leanmeta{\LH{CC2}}
\end{proposition}

\begin{proof}
The denominator $|B|$ is shared. Thus $\mathrm{ADQ}(I) < \mathrm{ADQ}(J)$ is equivalent to a strict inequality between the corresponding achievable-behavior set cardinalities.
\end{proof}

\subsection{Corollaries for Practice}

\begin{corollary}[Cardinality Criterion for Exact Minimization]\label{cor:practice-diagnostic}
\claimstamp{C\ref{cor:practice-diagnostic}}{H=succinct-hard}
In the mechanized Boolean-coordinate model, existence of a sufficient set of size at most $k$ is equivalent to the relevance set having cardinality at most $k$.
\leanmeta{\LH{S2P2}}
\end{corollary}

\begin{proof}
By \LH{S2P2}, sufficiency of size $\le k$ is equivalent to a relevance-cardinality bound $\le k$ in the Boolean-coordinate model.
\end{proof}

\begin{corollary}[Bounded-Regime Tractability]\label{cor:practice-bounded}
\claimstamp{C\ref{cor:practice-bounded}}{H=tractable-structured}
When the bounded-action or explicit-state conditions of Theorem~\ref{thm:tractable} hold, minimal modeling can be solved in polynomial time in the stated input size.
\leanmeta{\LH{CC81}}
\end{corollary}

\begin{proof}
This is the bounded-regime branch of Theorem~\ref{thm:tractable}, mechanized as \LH{CC81}.
\end{proof}

\begin{corollary}[Separable-Utility Tractability]\label{cor:practice-structured}
\claimstamp{C\ref{cor:practice-structured}}{H=tractable-structured}
When utility is separable with explicit factors, sufficiency checking is polynomial in the explicit-state regime.
\leanmeta{\LH{CC82}}
\end{corollary}

\begin{proof}
This is the separable-utility branch of Theorem~\ref{thm:tractable}, mechanized as \LH{CC82}.
\end{proof}

\begin{corollary}[Tree-Structured Tractability]\label{cor:practice-tree}
\claimstamp{C\ref{cor:practice-tree}}{H=tractable-structured}
When utility factors form a tree structure with explicit local factors, sufficiency checking is polynomial in the explicit-state regime.
\leanmeta{\LH{CC84}}
\end{corollary}

\begin{proof}
This is the tree-factor branch of Theorem~\ref{thm:tractable}, mechanized as \LH{CC84}.
\end{proof}

\begin{corollary}[Low-Tensor-Rank Tractability]\label{cor:practice-tensor}
\claimstamp{C\ref{cor:practice-tensor}}{H=tractable-structured}
When utility has tensor rank $R$ with explicit factor decomposition $U(a,s) = \sum_{r=1}^R w_r f_r(a) \prod_{i=1}^n g_{ri}(s_i)$, sufficiency checking is polynomial in $|A| \cdot R \cdot n$.
\leanmeta{\LH{CC71}}
\end{corollary}

\begin{proof}
This is the low-tensor-rank branch of Theorem~\ref{thm:tractable}, mechanized as \LH{CC71}.
\end{proof}

\begin{corollary}[Bounded-Treewidth Tractability]\label{cor:practice-treewidth}
\claimstamp{C\ref{cor:practice-treewidth}}{H=tractable-structured}
When utility decomposes as pairwise interactions on a graph of treewidth $w$, sufficiency checking is polynomial in $n \cdot k^{w+1}$.
\leanmeta{\LH{CC73},\LH{CC75}}
\end{corollary}

\begin{proof}
This is the bounded-treewidth branch of Theorem~\ref{thm:tractable}, mechanized as \LH{CC73}.
\end{proof}

\begin{corollary}[Coordinate-Symmetry Tractability]\label{cor:practice-symmetry}
\claimstamp{C\ref{cor:practice-symmetry}}{H=tractable-structured}
When state space is $S = [k]^d$ with coordinate-permutation-invariant utility, sufficiency checking is polynomial in $\binom{d+k-1}{k-1}^2 \cdot |A|$, which is polynomial in $d$ for fixed $k$.
\leanmeta{\LH{CC72}}
\end{corollary}

\begin{proof}
This is the coordinate-symmetry branch of Theorem~\ref{thm:tractable}, mechanized as \LH{CC72}.
\end{proof}

\begin{corollary}[Mechanized Hard Family]\label{cor:practice-unstructured}
There is a machine-checked family of reduction instances where, for non-tautological source formulas, every coordinate is relevant ($k^*=n$), exhibiting worst-case boundary complexity.
\leanmeta{\LH{CC26}}
\end{corollary}

\begin{proof}
The strengthened reduction proves that non-tautological source formulas induce instances where every coordinate is relevant; this is mechanized as \LH{CC26}.
\end{proof}

\subsection{Hardness Distribution: Right Place vs Wrong Place}\label{sec:hardness-distribution}

\begin{definition}[Hardness Distribution]\label{def:hardness-distribution}
Let $P$ be a problem family under the succinct encoding of Section~\ref{sec:encoding}. In this section, baseline hardness $H(P;n)$ denotes worst-case computational step complexity on instances with $n$ coordinates (equivalently, as a function of succinct input length $L$) in the fixed encoding regime. A \emph{solution architecture} $S$ partitions this baseline hardness into:
\begin{itemize}
\item $H_{\text{central}}(S)$: hardness paid once, at design time or in a shared component
\item $H_{\text{distributed}}(S)$: hardness paid per use site
\end{itemize}
For $n$ use sites, total realized hardness is:
\[H_{\text{total}}(S) = H_{\text{central}}(S) + n \cdot H_{\text{distributed}}(S)\]
\end{definition}

\begin{proposition}[Baseline Lower-Bound Principle]\label{prop:hardness-conservation}
\claimstamp{P\ref{prop:hardness-conservation}}{H=cost-growth}
For any problem family $P$ measured by $H(P;n)$ above, any solution architecture $S$ and any number of use sites $n \ge 1$, if $H_{\text{total}}(S)$ is measured in the same worst-case step units over the same input family, then:
\[
H_{\text{total}}(S) = H_{\text{central}}(S) + n \cdot H_{\text{distributed}}(S) \geq H(P;n).
\]
For SUFFICIENCY-CHECK, Theorem~\ref{thm:dichotomy} provides the baseline on the hard succinct family: $H(\text{SUFFICIENCY-CHECK};n)=2^{\Omega(n)}$ under ETH.
\leanmeta{\LH{HD23},\LH{HD25}}
\end{proposition}

\begin{proof}
By definition, $H(P;n)$ is a worst-case lower bound for correct solutions in this encoding regime and cost metric. Any such solution architecture decomposes total realized work as $H_{\text{central}} + n\cdot H_{\text{distributed}}$, so that total cannot fall below the baseline.
\end{proof}

\begin{definition}[Hardness Efficiency]\label{def:hardness-efficiency}
The \emph{hardness efficiency} of solution $S$ with $n$ use sites is:
\[\eta(S, n) = \frac{H_{\text{central}}(S)}{H_{\text{central}}(S) + n \cdot H_{\text{distributed}}(S)}\]
\leanmeta{\LH{HD11}}
\end{definition}

\begin{proposition}[Efficiency Equivalence]\label{prop:hardness-efficiency-interpretation}
\claimstamp{P\ref{prop:hardness-efficiency-interpretation}}{H=cost-growth}
For fixed $n$ and positive total hardness, larger $\eta(S,n)$ is equivalent to a larger central share of realized hardness.
\leanmeta{\LH{HD11}}
\end{proposition}

\begin{proof}
From Definition~\ref{def:hardness-efficiency}, $\eta(S,n)$ is exactly the fraction of total realized hardness paid centrally.
\end{proof}

\begin{definition}[Right vs Wrong Hardness Placement]\label{def:right-wrong-placement}
For a solution architecture $S$ in this linear model:
\[
\text{right hardness} \iff H_{\mathrm{distributed}}(S)=0,\qquad
\text{wrong hardness} \iff H_{\mathrm{distributed}}(S)>0.
\]
\leanmeta{\LH{HD12},
\LH{HD13}}
\end{definition}

\begin{theorem}[Centralization Dominance]\label{thm:centralization-dominance}
Let $S_{\mathrm{right}}, S_{\mathrm{wrong}}$ be architectures over the same problem family with
\[
H_{\mathrm{distributed}}(S_{\mathrm{right}})=0,\quad
H_{\mathrm{central}}(S_{\mathrm{right}})>0,\quad
H_{\mathrm{distributed}}(S_{\mathrm{wrong}})>0,
\]
and let $n > \max\!\bigl(1, H_{\mathrm{central}}(S_{\mathrm{right}})\bigr)$. Then:
\begin{enumerate}
\item Lower total realized hardness:
\[
H_{\text{total}}(S_{\mathrm{right}}) < H_{\text{total}}(S_{\mathrm{wrong}})
\]
\item Fewer error sites: errors in centralized components affect 1 location; errors in distributed components affect $n$ locations
\item Quantified leverage: moving one unit of work from distributed to central saves exactly $n-1$ units of total realized hardness
\end{enumerate}
\leanmeta{\LH{HD1},
\LH{HD2}}
\end{theorem}

\begin{proof}
(1) and (2) are exactly the bundled theorem
\LH{HD1}.
(3) is exactly
\LH{HD2}.
\end{proof}

\begin{corollary}[Right-Place vs Wrong-Place Hardness]\label{cor:right-wrong-hardness}
\claimstamp{C\ref{cor:right-wrong-hardness}}{H=cost-growth}
In the linear model, a right-hardness architecture strictly dominates a wrong-hardness architecture once use-site count exceeds central one-time hardness. Formally, for architectures $S_{\mathrm{right}}, S_{\mathrm{wrong}}$ over the same problem family, if $S_{\mathrm{right}}$ has right hardness, $S_{\mathrm{wrong}}$ has wrong hardness, and $n > H_{\mathrm{central}}(S_{\mathrm{right}})$, then
\[
H_{\mathrm{central}}(S_{\mathrm{right}}) + n\,H_{\mathrm{distributed}}(S_{\mathrm{right}})
<
H_{\mathrm{central}}(S_{\mathrm{wrong}}) + n\,H_{\mathrm{distributed}}(S_{\mathrm{wrong}}).
\]
\leanmeta{\LH{HD19}}
\end{corollary}

\begin{proof}
This is the mechanized theorem \LH{HD19}.
\end{proof}

\begin{proposition}[Dominance Modes]\label{prop:dominance-modes}
\claimstamp{P\ref{prop:dominance-modes}}{H=cost-growth}
This section uses two linear-model dominance modes plus generalized nonlinear dominance and boundary modes:
\begin{enumerate}
\item \textbf{Strict threshold dominance:} Corollary~\ref{cor:right-wrong-hardness} gives strict inequality once $n > H_{\mathrm{central}}(S_{\mathrm{right}})$.
\item \textbf{Global weak dominance:} under the decomposition identity used in \LH{HD3}, centralized hardness placement is never worse for all $n\ge 1$.
\item \textbf{Generalized nonlinear dominance:} under bounded-vs-growing site-cost assumptions (Theorem~\ref{thm:generalized-dominance}), right placement strictly dominates beyond a finite threshold without assuming linear per-site cost.
\item \textbf{Generalized boundary mode:} without those growth-separation assumptions, strict dominance is not guaranteed (Proposition~\ref{prop:generalized-assumption-boundary}).
\end{enumerate}
\end{proposition}

\begin{proof}
Part (1) is Corollary~\ref{cor:right-wrong-hardness}. Part (2) is exactly \LH{HD3}. Part (3) is Theorem~\ref{thm:generalized-dominance}. Part (4) is Proposition~\ref{prop:generalized-assumption-boundary}.
\end{proof}

\textbf{Illustrative Instantiation (Type Systems).} Consider a capability $C$ (e.g., provenance tracking) with one-time central cost $H_{\text{central}}$ and per-site manual cost $H_{\text{distributed}}$:

\begin{center}
\begin{tabular}{lcc}
\toprule
\textbf{Approach} & $H_{\text{central}}$ & $H_{\text{distributed}}$ \\
\midrule
Native type system support & High (learning cost) & Low (type checker enforces) \\
Manual implementation & Low (no new concepts) & High (reimplement per site) \\
\bottomrule
\end{tabular}
\end{center}
The table is schematic; the formal statement is Corollary~\ref{cor:type-system-threshold}. The type-theoretic intuition behind this instantiation is consistent with classic accounts of abstraction and static interface guarantees \cite{Cardelli1985, reynolds1983types, pierce2002types, liskov1994behavioral, siek2006gradual, gamma1994design}.

\begin{corollary}[Type-System Threshold]\label{cor:type-system-threshold}
\claimstamp{C\ref{cor:type-system-threshold}}{H=cost-growth}
For the formal native-vs-manual architecture instance, native support has lower total realized cost for all
\[
n > H_{\mathrm{baseline}}(P),
\]
where $H_{\mathrm{baseline}}(P)$ corresponds to \LH{HD15}.
\leanmeta{\LH{HD15}}
\end{corollary}

\begin{proof}
Immediate from \LH{HD15}.
\end{proof}

\subsection{Extension: Non-Additive Site-Cost Models}\label{sec:nonadditive-site-costs}

\begin{definition}[Generalized Site Accumulation]\label{def:generalized-site-accumulation}
Let $C_S : \mathbb{N} \to \mathbb{N}$ be a per-site accumulation function for architecture $S$. Define generalized total realized hardness by
\[
H_{\text{total}}^{\mathrm{gen}}(S,n) = H_{\text{central}}(S) + C_S(n).
\]
\end{definition}

\begin{definition}[Eventual Saturation]\label{def:eventual-saturation}
A cost function $f : \mathbb{N}\to\mathbb{N}$ is \emph{eventually saturating} if there exists $N$ such that for all $n\ge N$, $f(n)=f(N)$.
\end{definition}

\begin{theorem}[Generalized Dominance by Growth Separation]\label{thm:generalized-dominance}
\claimstamp{T\ref{thm:generalized-dominance}}{H=cost-growth}
Let
\[
H_{\text{total}}^{\mathrm{gen}}(S,n)=H_{\text{central}}(S)+C_S(n).
\]
For two architectures $S_{\mathrm{right}},S_{\mathrm{wrong}}$, suppose there exists $B\in\mathbb{N}$ such that:
\begin{enumerate}
\item $C_{S_{\mathrm{right}}}(m)\le B$ for all $m$ (bounded right-side per-site accumulation),
\item $m \le C_{S_{\mathrm{wrong}}}(m)$ for all $m$ (identity-lower-bounded wrong-side growth).
\end{enumerate}
Then for every
\[
n > H_{\text{central}}(S_{\mathrm{right}})+B,
\]
one has
\[
H_{\text{total}}^{\mathrm{gen}}(S_{\mathrm{right}},n)
<
H_{\text{total}}^{\mathrm{gen}}(S_{\mathrm{wrong}},n).
\]
\leanmeta{\LH{HD9}}
\end{theorem}

\begin{proof}
This is exactly the mechanized theorem
\LH{HD9}.
\end{proof}

\begin{corollary}[Eventual Generalized Dominance]\label{cor:generalized-eventual-dominance}
\claimstamp{C\ref{cor:generalized-eventual-dominance}}{H=cost-growth}
If condition (1) above holds and there exists $N$ such that condition (2) holds for all $m\ge N$, then there exists $N_0$ such that for all $n\ge N_0$,
\[
H_{\text{total}}^{\mathrm{gen}}(S_{\mathrm{right}},n)
<
H_{\text{total}}^{\mathrm{gen}}(S_{\mathrm{wrong}},n).
\]
\leanmeta{\LH{HD10}}
\end{corollary}

\begin{proof}
Immediate from
\LH{HD10}.
\end{proof}

\begin{proposition}[Generalized Assumption Boundary]\label{prop:generalized-assumption-boundary}
\claimstamp{P\ref{prop:generalized-assumption-boundary}}{H=cost-growth}
In the generalized model, strict right-vs-wrong dominance is not unconditional. There are explicit counterexamples:
\begin{enumerate}
\item If wrong-side growth lower bounds are dropped, right-side strict dominance can fail for all $n$.
\item If right-side boundedness is dropped, strict dominance can fail for all $n$ even when wrong-side growth is linear.
\end{enumerate}
\leanmeta{\LH{HD8},
\LH{HD7}}
\end{proposition}

\begin{proof}
This is exactly the pair of mechanized boundary theorems listed above.
\end{proof}

\begin{theorem}[Linear Model: Saturation iff Zero Distributed Hardness]\label{thm:linear-saturation-iff-zero}
\claimstamp{T\ref{thm:linear-saturation-iff-zero}}{H=cost-growth}
In the linear model of this section,
\[
H_{\text{total}}(S,n)=H_{\text{central}}(S)+n\cdot H_{\text{distributed}}(S),
\]
the function $n\mapsto H_{\text{total}}(S,n)$ is eventually saturating if and only if $H_{\text{distributed}}(S)=0$.
\leanmeta{\LH{HD19}}
\end{theorem}

\begin{proof}
This is exactly the mechanized equivalence theorem above.
\end{proof}

\begin{theorem}[Generalized Model: Saturation is Possible]\label{thm:generalized-saturation-possible}
\claimstamp{T\ref{thm:generalized-saturation-possible}}{H=cost-growth}
There exists a generalized site-cost model with eventual saturation. In particular, for
\[
C_K(n)=\begin{cases}
n, & n\le K\\
K, & n>K,
\end{cases}
\]
both $C_K$ and $n\mapsto H_{\text{central}}+C_K(n)$ are eventually saturating.
\leanmeta{\LH{HD20},
\LH{HD6}}
\end{theorem}

\begin{proof}
This is the explicit construction mechanized in Lean.
\end{proof}

\begin{corollary}[Positive Linear Slope Cannot Represent Saturation]\label{cor:linear-positive-no-saturation}
\claimstamp{C\ref{cor:linear-positive-no-saturation}}{H=cost-growth}
No positive-slope linear per-site model can represent the saturating family above for all $n$.
\leanmeta{\LH{HD16}}
\end{corollary}

\begin{proof}
This follows from the mechanized theorem that any linear representation of the saturating family must have zero slope.
\end{proof}

\paragraph{Mechanized strengthening reference.}
The strengthened all-coordinates-relevant reduction is presented in Section~\ref{sec:hardness} (``Mechanized strengthening'') and formalized in \texttt{Reduction\_AllCoords.lean} via \texttt{all\_coords\_relevant\_of\_not\_tautology}.

The next section develops the major practical consequence of this framework: the Simplicity Tax Theorem.
